% SOURCE_AUTHORITY: sga5_fr_workpass.tex, lines 11771--11797 (LF-normalized unit).
% SOURCE_SHA256: 791F4EFFC5E02832D5D77ED03518C8156D6F07E4C8238B03545DB93D883FBB28
\emph{Demostración.} Es claro que existe un homomorfismo natural
\[
\operatorname{Coker}(K(i))\xrightarrow{\alpha}K(\cA/\cB).
\]
Por otra parte, existe una aplicación evidente
\[
\Ob(\cA/\cB)\xrightarrow{\gamma}\operatorname{Coker}(K(i)),
\]
puesto que $\Ob(\cA/\cB)=\Ob(\cA)$. Sean $X,Y\in\Ob(\cA/\cB)$, y supongamos que son isomorfos. Queremos demostrar que $\gamma(X)=\gamma(Y)$. La categoría $\cA/\cB$ se obtiene mediante un cálculo de fracciones por la derecha o por la izquierda, siendo el sistema multiplicativo $S$ el conjunto de los morfismos $f$ que están contenidos en un triángulo distinguido $(X',Y',Z',f,g,h)$, donde $Z'$ es un objeto de $\cB$. ($S$ es un sistema multiplicativo saturado.) Resulta entonces que existen $X_1$ y
\[
X_1\xrightarrow{s}X,\qquad X_1\xrightarrow{t}Y,
\]
tales que $s,t\in S$. Basta, pues, demostrar que $\gamma(X)=\gamma(X_1)$, lo que resulta del hecho de que existe un triángulo distinguido de la forma
\[
(X_1,X,Z,s,\ldots),
\]
con $Z\in\Ob\cB$.

La aplicación $\gamma$ es una función aditiva. En efecto, siendo $(X,Y,Z,f,g,h)$ un triángulo distinguido de $\cA/\cB$, hay que demostrar que
\[
\gamma(X)+\gamma(Z)=\gamma(Y).
\]
Pero la relación se sigue ahora porque el triángulo $(X,Y,Z,f,g,h)$ es isomorfo a un triángulo distinguido de $\cA$. Por consiguiente, se obtiene
\[
\beta:K(\cA/\cB)\longrightarrow \operatorname{Coker}(K(i)),
\]
y es inmediato comprobar que $\beta$ es inversa de $\alpha$.
